\documentclass[11pt,a4paper]{article}
\usepackage[T1]{fontenc}
\usepackage[utf8]{inputenc}
\usepackage[main=italian,provide=*]{babel}
\usepackage{amsmath,amssymb}
\usepackage{geometry}
\geometry{margin=2.5cm}
\usepackage{booktabs}
\usepackage{seqsplit}
\usepackage{hyperref}
\hypersetup{colorlinks=true,linkcolor=blue,citecolor=blue,urlcolor=blue}

\title{Manuale d'uso --- correlazioni dinamiche trimero anello\\
\large versione a stato esatto (\texttt{prepare\_state})}
\author{Note di lavoro --- Tesi triennale, Samuele Schiavi\\ Relatore: Prof.\ Alessandro Chiesa}
\date{}

\begin{document}
\maketitle

\section*{Scopo}

Questo manuale spiega come eseguire la simulazione delle correlazioni dinamiche
$C_{ij}^{\alpha\beta}(t)$ per il trimero ad anello nella sua versione
\emph{originale}: lo stato fondamentale $\ket{\psi_0}$ è preparato sul registro
con le ampiezze esatte (\texttt{QuantumCircuit.prepare\_state}), non con un
ansatz VQE. È la versione da usare quando interessa isolare l'errore di
Trotter e quello statistico (shot), senza mescolarli con un residuo di
fedeltà dell'ansatz variazionale. Per la versione con il circuito VQE reale
(\texttt{W-2q.6}), vedi il manuale gemello
\texttt{\seqsplit{manuale\_uso\_correlazioni\_trimero\_anello\_vqe.tex}}.

\section*{Prerequisiti}

\textbf{File necessari nella stessa cartella} (tutti nella fase
\emph{correlazioni dinamiche}, trimero anello):
\begin{itemize}
\item \texttt{trimer\_ring\_exact.py} --- Hamiltoniana \texttt{trimer\_hamiltonian\_dm}.
\item \texttt{trotter\_trimero\_anello.py} --- evoluzione $U(t)=e^{-iHt}$ via Suzuki--Trotter.
\item \texttt{circuito\_correlazioni\_trimero\_anello.py} --- il circuito Hadamard test a 4 qubit.
\item \texttt{validate\_circuito\_correlazioni.py} --- riferimento classico esatto (\texttt{scipy.linalg.expm}).
\item \texttt{validate\_shot\_noise\_trimero\_anello.py} --- validazione a shot finiti.
\item \texttt{scan81\_trimero\_anello.py} --- scan completo delle 81 combinazioni.
\item \texttt{generate\_figure\_shotnoise.py}, \texttt{generate\_figure\_scan81.py} --- figure.
\end{itemize}

\textbf{Librerie Python}: \texttt{qiskit} (2.x), \texttt{qiskit-aer}, \texttt{numpy},
\texttt{scipy}, \texttt{matplotlib} (solo per le figure). Verificare con:
\begin{verbatim}
python -c "import qiskit; print(qiskit.__version__)"
\end{verbatim}

\textbf{Punto di lavoro di riferimento} (usato ovunque in questa fase, salvo
diversa indicazione): $J=1$, $J'=0.4$, $b=b_c=2.4$, $D=0.15$ (DM, Opzione B),
$t=1.3$, $N=200$ passi di Trotter.

\section*{Passo 1 --- un correlatore singolo, da riga di comando Python}

Per calcolare un correlatore specifico, ad esempio $C_{11}^{yy}(t{=}1.3)$:
\begin{verbatim}
from circuito_correlazioni_trimero_anello import ground_state, correlator_from_circuit

J, Jp, b, D = 1.0, 0.4, 2.4, 0.15
psi0, E = ground_state(J, Jp, b, D)

C = correlator_from_circuit(1, 'y', 1, 'y', t=1.3, N=200, J=J, Jp=Jp, b=b, D=D, psi0=psi0)
print(C)   # atteso: un numero complesso, |C| dell'ordine di alcuni decimi
\end{verbatim}
Non passare \texttt{ansatz\_params}: di default resta \texttt{None} e la
preparazione avviene con \texttt{prepare\_state(psi0)}, ampiezze esatte.

\section*{Passo 2 --- validazione contro il riferimento classico esatto}

\begin{verbatim}
python validate_circuito_correlazioni.py
\end{verbatim}
Stampa un confronto, per sei correlatori rappresentativi e tre valori di $t$,
tra il valore classico esatto ($\texttt{scipy.linalg.expm}$, nessun errore di
Trotter) e quello dal circuito ($N=200$): il residuo atteso è dell'ordine di
$10^{-5}$--$10^{-6}$ (errore di Trotter puro). In coda, uno studio di
convergenza in $N$ su $C_{11}^{yy}(t{=}1.3)$: l'errore deve dimezzarsi
raddoppiando $N$ (Trotter al primo ordine, $O(1/N)$ sull'ampiezza).

\section*{Passo 3 --- validazione a shot finiti}

\begin{verbatim}
python validate_shot_noise_trimero_anello.py
\end{verbatim}
Separa i due contributi di errore: quello algoritmico (Trotter, statevector
vs esatto classico) e quello statistico (shot vs statevector, $8192$ shots).
Al punto di lavoro standard l'errore di Trotter domina nettamente su quello
statistico (rapporto $\sim16\times$, vedi
\texttt{\seqsplit{circuito\_correlazioni\_trimero\_anello\_spiegato.tex}}) --
regime opposto a quello del dimero.

\section*{Passo 4 --- scan completo delle 81 combinazioni}

\begin{verbatim}
python scan81_trimero_anello.py
\end{verbatim}
Richiede qualche minuto (81 combinazioni $\times$ 2 parti Re/Im $\times$
(statevector + 8192 shots), con cache di trascompilazione). Stampa un
riepilogo (zeri strutturali confermati, errori medi/massimi, correlatori più
"ricchi") e salva due file nella cartella corrente:
\begin{itemize}
\item \texttt{scan81\_results.npz} --- matrici $9\times9$ (una riga/colonna per
ogni coppia sito-componente) di riferimento classico, statevector e shot,
più le statistiche di errore. Si carica con \texttt{numpy.load(...)}.
\item \texttt{scan81\_results.json} --- stessi dati, riga per riga, in formato
leggibile/ispezionabile a mano.
\end{itemize}
\textbf{Attenzione}: lo script salva con percorso \emph{relativo} --- va
eseguito da dentro la cartella dove si vogliono i risultati (non da un'altra
posizione), altrimenti Python solleva \texttt{FileNotFoundError}.

\section*{Passo 5 --- figure}

\begin{verbatim}
python generate_figure_shotnoise.py
python generate_figure_scan81.py
\end{verbatim}
Il secondo richiede che \texttt{scan81\_results.npz} sia già presente (Passo 4
completato). Producono i PNG usati nei documenti di teoria/spiegazione.

\section*{Problemi comuni}

\begin{itemize}
\item \texttt{ModuleNotFoundError} su uno dei moduli locali --- verificare di
eseguire dalla cartella che contiene tutti i file elencati in
"Prerequisiti", non da una sottocartella o da altrove.
\item \texttt{FileNotFoundError} su \texttt{scan81\_results.npz} in uno script
successivo --- va prima generato eseguendo \texttt{scan81\_trimero\_anello.py}
(Passo 4); non è incluso di default perché il calcolo richiede qualche minuto.
\item Risultati leggermente diversi da quelli riportati nei documenti di teoria
--- controllare punto di lavoro ($J,J',b,D$) e $t,N$: sono hard-coded in testa
a ogni script, non parametri da riga di comando.
\end{itemize}

\section*{Riferimenti}

Teoria e derivazioni: \texttt{\seqsplit{simmetrie\_correlatori\_trimero\_anello.tex}},
\texttt{\seqsplit{circuito\_correlazioni\_trimero\_anello\_spiegato.tex}}. Catalogo
completo della fase: \texttt{\seqsplit{trimero\_anello\_correlazioni.tex}}.
Cronologia delle decisioni: \texttt{log\_decisioni.md}.

\end{document}